% =====================================================================
%  AEC Scholar — generic academic article template
%  Replace the generic `article` class with your target journal's class
%  (elsarticle / IEEEtran / ascelike-new / sn-jnl) before submission.
%  Compile:  pdflatex -> bibtex -> pdflatex -> pdflatex   (or: latexmk -pdf)
% =====================================================================
\documentclass[11pt,a4paper]{article}

% --- Encoding & fonts (for pdfLaTeX; switch to fontspec under Lua/XeLaTeX) ---
\usepackage[utf8]{inputenc}
\usepackage[T1]{fontenc}
\usepackage{lmodern}
\usepackage{microtype}                 % better justification

% --- Layout ---
\usepackage[margin=2.5cm]{geometry}
\usepackage{setspace}                  % \onehalfspacing for review copies
\usepackage{lineno}                    % \linenumbers for review (many journals require)

% --- Graphics, tables, math, units ---
\usepackage{graphicx}
\graphicspath{{figs/}}
\usepackage{booktabs}                  % professional tables (no vertical rules)
\usepackage{subcaption}                % subfigures
\usepackage{amsmath,amssymb}
\usepackage{siunitx}                   % \qty{}{}, \num{}, S-columns
\sisetup{detect-all=true}

% --- References & links (load hyperref before cleveref) ---
\usepackage[numbers,sort&compress]{natbib}   % or author-year per venue
\usepackage[hidelinks]{hyperref}
\usepackage[capitalise,nameinlink]{cleveref}

% --- Handy review note macro (remove before submission) ---
\usepackage{xcolor}
\newcommand{\todo}[1]{\textcolor{red}{[\textbf{TODO:} #1]}}

% =====================================================================
\title{A Concise, Specific, Keyword-Rich Title:\\ Method and Context Make It Discoverable}
\author{First Author\thanks{Corresponding author: author@coera.io} \and Second Author}
\date{\today}

\begin{document}
\maketitle
% \linenumbers   % <- enable for review submission

\begin{abstract}
\noindent
% Context/problem -> gap/objective -> method -> key results (concrete) -> contribution.
State the problem and why it matters, the specific gap, the method, the key quantitative findings,
and the contribution — in 150--250 words, self-contained, no citations, no undefined acronyms.
\end{abstract}

\noindent\textbf{Keywords:} keyword one; keyword two; keyword three; keyword four; keyword five

% =====================================================================
\section{Introduction}\label{sec:intro}
% CARS move: (1) establish the territory, (2) establish the niche / GAP, (3) occupy it.
Context and significance of the topic for the built environment.\ \todo{cite}
The specific gap or unresolved problem in prior work.
This paper contributes \todo{state the explicit contribution}.
The remainder is organised as follows: \Cref{sec:related} reviews related work; \Cref{sec:method}
describes the methodology; \Cref{sec:results} presents results; \Cref{sec:discussion} discusses them;
and \Cref{sec:conclusion} concludes.

\section{Related Work / Background}\label{sec:related}
Synthesise prior work thematically and position the gap. Define the conceptual framework.

\section{Methodology}\label{sec:method}
Describe the research design with enough detail to reproduce: data/sample, tools (with versions),
procedure, analysis, and validity/reliability measures. Reference the workflow in \Cref{fig:workflow}.

\begin{figure}[t]
  \centering
  % \includegraphics[width=0.8\linewidth]{workflow.pdf}
  \fbox{\parbox[c][4cm][c]{0.8\linewidth}{\centering\itshape Replace with research-workflow figure (vector PDF).}}
  \caption{Overview of the research workflow.}
  \label{fig:workflow}
\end{figure}

\section{Results}\label{sec:results}
Report findings without interpretation. Lead with the answer to each research question.

\begin{table}[t]
  \centering
  \caption{Summary of results.}
  \label{tab:results}
  \begin{tabular}{l S[table-format=3.1] S[table-format=2.2]}
    \toprule
    Case & {Metric A (\unit{\kilo\watt\hour\per\meter\squared})} & {Metric B} \\
    \midrule
    Baseline & 145.0 & 1.00 \\
    Proposed & 118.4 & 0.82 \\
    \bottomrule
  \end{tabular}
\end{table}

\section{Discussion}\label{sec:discussion}
Interpret the results, compare with prior work, state theoretical and practical implications, and be
explicit about limitations (especially generalisability) and future work.

\section{Conclusion}\label{sec:conclusion}
Restate the contribution and key findings, implications, limitations, and future directions — no new data.

% --- Declarations (many venues now require these) ---
\section*{Declarations}
\textbf{Funding:} \todo{sources + grant numbers}.\\
\textbf{Conflicts of interest:} The authors declare none.\\
\textbf{Data availability:} \todo{repository + DOI, or reason data cannot be shared}.\\
\textbf{CRediT author statement:} \todo{roles per author}.\\
\textbf{Use of AI tools:} \todo{disclose per venue policy, if applicable}.

\section*{Acknowledgements}
Non-author contributions and support.

% --- Bibliography ---
\bibliographystyle{plainnat}   % swap for the venue style (elsarticle-num / IEEEtran / ascelike)
\bibliography{references}      % references.bib in the same folder

\end{document}
